\begin{thebibliography}{99}

\bibitem{black-scholes}
F. Black, M. Scholes.
The pricing of options and corporate liabilities.
\emph{Journal of Political Economy}, 81 (3), 637--654, 1973.

\bibitem{merton1973}
R. C. Merton.
Theory of rational option pricing.
\emph{Bell Journal of Economics and Management Science}, 4 (1),
141--183, 1973.

\bibitem{baaquie2004}
B. E. Baaquie.
\emph{Quantum Finance: Path Integrals and Hamiltonians for Options
and Interest Rates}.
Cambridge University Press, Cambridge, 2004.

\bibitem{linetsky1998}
V. Linetsky.
The path integral approach to financial modeling and option pricing.
\emph{Computational Economics}, 11, 129--163, 1998.

\bibitem{linetsky2004}
V. Linetsky.
Spectral expansions for Asian (average price) options.
\emph{Operations Research}, 52 (6), 856--867, 2004.

\bibitem{kemna-vorst}
A. G. Z. Kemna, A. C. F. Vorst.
A pricing method for options based on average asset values.
\emph{Journal of Banking and Finance}, 14 (1), 113--129, 1990.

\bibitem{geman-yor}
H. Geman, M. Yor.
Bessel processes, Asian options, and perpetuities.
\emph{Mathematical Finance}, 3 (4), 349--375, 1993.

\bibitem{kac1949}
M. Kac.
On distributions of certain Wiener functionals.
\emph{Transactions of the American Mathematical Society}, 65 (1),
1--13, 1949.

\bibitem{feynman-hibbs}
R. P. Feynman, A. R. Hibbs.
\emph{Quantum Mechanics and Path Integrals}.
McGraw-Hill, Nueva York, 1965.

\bibitem{shreve}
S. E. Shreve.
\emph{Stochastic Calculus for Finance II: Continuous-Time Models}.
Springer Finance. Springer, Nueva York, 2004.

\bibitem{hull}
J. C. Hull.
\emph{Options, Futures, and Other Derivatives}.
10\textsuperscript{a} edición, Pearson, Boston, 2017.

\bibitem{glasserman}
P. Glasserman.
\emph{Monte Carlo Methods in Financial Engineering}.
Springer, Nueva York, 2003.

\bibitem{broadie-glasserman-kou}
M. Broadie, P. Glasserman, S. Kou.
A continuity correction for discrete barrier options.
\emph{Mathematical Finance}, 7 (4), 325--349, 1997.

\bibitem{tfgm}
R. Cepeda Rocandio.
\emph{Valoración de opciones europeas: del modelo binomial al modelo
de Black--Scholes}.
Trabajo Fin de Grado en Matemáticas, Universidad de Oviedo, 2027.

\end{thebibliography}
